
% BEGIN bib
 
\par\medskip\noindent\textbf{References}\par
  
\par\noindent  [1]  Filippo Bonchi, Cipriano Junior Cioffo, Alessandro Di Giorgio, and Elena Di Lavore.  Tape Diagrams for Monoidal Monads.  In Corina Cîrstea and Alexander Knapp, editors, 11th Conference on Algebra and Coalgebra in Computer Science (CALCO 2025), volume 342 of Leibniz International Proceedings in Informatics (LIPIcs), pages 11:1–11:24, Dagstuhl, Germany, 2025. Schloss Dagstuhl – Leibniz-Zentrum für Informatik.  doi:10.4230/LIPIcs.CALCO.2025.11.  
 
\par\noindent  [2]  Anne Broadbent and Martti Karvonen.  Categorical composable cryptography: extended version.  Logical Methods in Computer Science, Volume 19, Issue 4, December 2023.  doi:10.46298/lmcs-19(4:30)2023.  
 
\par\noindent  [3]  Ran Canetti.  Universally composable security: A new paradigm for cryptographic protocols.  In Proceedings 42nd IEEE Symposium on Foundations of Computer Science, pages 136–145. IEEE, 2001.  doi:10.1109/SFCS.2001.959888.  
 
\par\noindent  [4]  Titouan Carette, Dominic Horsman, and Simon Perdrix.  Szx-calculus: Scalable graphical quantum reasoning.  In MFCS 2019-44th International Symposium on Mathematical Foundations of Computer Science, volume 138, pages 55–1, 2019.  
 
\par\noindent  [5]  Apiwat Chantawibul and Pawel Sobocinski.  Monoidal multiplexing.  In B. Fischer and T. Uustalu, editors, Theoretical Aspects of Computing ? ICTAC 2018, volume 11187, pages 116–131. Springer, October 2018.  doi:10.1007/978-3-030-02508-3\_7.  
 
\par\noindent  [6]  Bob Coecke, Tobias Fritz, and Robert W Spekkens.  A mathematical theory of resources.  Information and Computation, 250:59–86, 2016.  doi:10.1016/J.IC.2016.02.008.  
 
\par\noindent  [7]  Cole Comfort, Antonin Delpeuch, and Jules Hedges.  Sheet diagrams for bimonoidal categories.  arXiv preprint arXiv:2010.13361, 2020.  
 
\par\noindent  [8]  Ugo Dal Lago, Furio Honsell, Marina Lenisa, and Paolo Pistone.  On Quantitative Algebraic Higher-Order Theories.  In Amy P. Felty, editor, 7th International Conference on Formal Structures for Computation and Deduction (FSCD 2022), volume 228 of Leibniz International Proceedings in Informatics (LIPIcs), pages 4:1–4:18, Dagstuhl, Germany, 2022. Schloss Dagstuhl – Leibniz-Zentrum für Informatik.  doi:10.4230/LIPIcs.FSCD.2022.4.  
 
\par\noindent  [9]  Elena Di Lavore, Giovanni de Felice, and Mario Román.  Monoidal streams for dataflow programming.  In Proceedings of the 37th Annual ACM/IEEE Symposium on Logic in Computer Science, LICS ’22, New York, NY, USA, 2022. Association for Computing Machinery.  doi:10.1145/3531130.3533365.  
 
\par\noindent  [10]  Tobias Fritz.  A synthetic approach to Markov kernels, conditional independence and theorems on sufficient statistics.  Advances in Mathematics, 370:107239, 2020.  arXiv:1908.07021.  
 
\par\noindent  [11]  Steven Holtzen, Guy Van den Broeck, and Todd Millstein.  Scaling exact inference for discrete probabilistic programs.  Proceedings of the ACM on Programming Languages, 4(OOPSLA):1–31, 2020.  doi:10.1145/3428208.  
 
\par\noindent  [12]  Nicholas Gauguin Houghton-Larsen.  A mathematical framework for causally structured dilations and its relation to quantum self-testing.  arXiv preprint arXiv:2103.02302, 2021.  arXiv:2103.02302.  
 
\par\noindent  [13]  André Joyal and Ross Street.  The geometry of tensor calculus, I.  Advances in Mathematics, 88(1):55–112, July 1991.  doi:10.1016/0001-8708(91)90003-P.  
 
\par\noindent  [14]  Jonathan Katz and Yehuda Lindell.  Introduction to Modern Cryptography, Second Edition.  Chapman \& Hall/CRC, 2nd edition, 2014.  
 
\par\noindent  [15]  A.N. Kolmogorov and S.V. Fomin.  Introductory Real Analysis.  Selected Russian publications in the mathematical sciences. Prentice-Hall, 1970.  URL: https://books.google.nl/books?id=9l4PAQAAMAAJ.  
 
\par\noindent  [16]  Ugo Dal Lago, Zeinab Galal, and Giulia Giusti.  On computational indistinguishability and logical relations.  In Programming Languages and Systems: 22nd Asian Symposium, APLAS 2024, Kyoto, Japan, October 22-24, 2024, Proceedings, pages 241–263, Berlin, Heidelberg, 2024. Springer-Verlag.  doi:10.1007/978-981-97-8943-6\_12.  
 
\par\noindent  [17]  Gabriele Lobbia, Wojciech Różowski, Ralph Sarkis, and Fabio Zanasi.  Quantitative Monoidal Algebra: Axiomatising Distance with String Diagrams.  In Paweł Gawrychowski, Filip Mazowiecki, and Michał Skrzypczak, editors, 50th International Symposium on Mathematical Foundations of Computer Science (MFCS 2025), volume 345 of Leibniz International Proceedings in Informatics (LIPIcs), pages 68:1–68:21, Dagstuhl, Germany, 2025. Schloss Dagstuhl – Leibniz-Zentrum für Informatik.  doi:10.4230/LIPIcs.MFCS.2025.68.  
 
\par\noindent  [18]  Saunders Mac Lane.  Categories for the Working Mathematician.  Springer, 1998.  
 
\par\noindent  [19]  Ueli Maurer.  Constructive cryptography–a new paradigm for security definitions and proofs.  In Joint Workshop on Theory of Security and Applications, pages 33–56. Springer, 2011.  doi:10.1007/978-3-642-27375-9\_3.  
 
\par\noindent  [20]  Ueli Maurer and Renato Renner.  Abstract cryptography.  In Bernard Chazelle, editor, Innovations in Computer Science - ICS 2011, Tsinghua University, Beijing, China, January 7-9, 2011. Proceedings, pages 1–21. Tsinghua University Press, 2011.  URL: http://conference.iiis.tsinghua.edu.cn/ICS2011/content/papers/14.html.  
 
\par\noindent  [21]  Robin Piedeleu, Mateo Torres-Ruiz, Alexandra Silva, and Fabio Zanasi.  A complete axiomatisation of equivalence for discrete probabilistic programming.  In Programming Languages and Systems: 34th European Symposium on Programming, ESOP 2025, Held as Part of the International Joint Conferences on Theory and Practice of Software, ETAPS 2025, Hamilton, ON, Canada, May 3–8, 2025, Proceedings, Part II, pages 202–229, Berlin, Heidelberg, 2025. Springer-Verlag.  doi:10.1007/978-3-031-91121-7\_9.  
 
\par\noindent  [22]  Robin Piedeleu and Fabio Zanasi.  An Introduction to String Diagrams for Computer Scientists.  Elements in Applied Category Theory. Cambridge University Press, 2025.  
 
\par\noindent  [23]  P. Selinger.  A Survey of Graphical Languages for Monoidal Categories, volume 813 of Lecture Notes in Physics, pages 289–355.  Springer, Berlin, Heidelberg, 2010.  doi:10.1007/978-3-642-12821-9\_4.  
 
\par\noindent  [24]  Mateo Torres-Ruiz, Robin Piedeleu, Alexandra Silva, and Fabio Zanasi.  On iteration in discrete probabilistic programming.  In 9th International Conference on Formal Structures for Computation and Deduction (FSCD 2024), volume 299 of Leibniz International Proceedings in Informatics (LIPIcs), pages 20:1–20:21. Schloss Dagstuhl – Leibniz-Zentrum für Informatik, 2024.  doi:10.4230/LIPIcs.FSCD.2024.20.  
  
% END bib
